% Book 3, Year 1, chapter 3 -- functions and the call stack.
% Provenance (never printed): sources/book3/BRIEFS.md, chapter 3.
\chapter{Functions and the Call Stack}\label{ch:b1:c-functions-and-stack}

The function \code{swap} below exchanges its two parameters, and prints them to prove it. Back in \code{main},
\code{x} is still 1 and \code{y} is still 2. The function did its work on copies, which disappeared when it
returned. Where those copies lived, how long they lived, and why a function that calls itself a million times can
crash a program, is the matter of this chapter: it opens the machine room under every function call.

\omcode{code/bachelor-1/03-c-functions-and-stack/listings/swap_value.c}{2}{18}{A swap that
swaps nothing}\label{lst:b1:c-functions-and-stack:swap}

\omoutput{code/bachelor-1/03-c-functions-and-stack/outputs/swap_value.txt}{Swapped inside, unchanged
outside.}\label{out:b1:c-functions-and-stack:swap}

\begin{recall}
A recursive function calls itself on a smaller input and stops at a base case; its correctness is proved by
induction on the size of the input. The school volume drew the calls of a recursive Python function as a stack of
frames, one per call still running. This chapter shows what those frames are made of in a compiled program, and
\cref{ch:b1:induction-and-termination} returns to the proofs.
\end{recall}

\section{Declaring and defining functions}\label{sec:b1:c-functions-and-stack:declaring}

A C function has a type for each parameter and a type for its result, and the \omterm{def:b1:compilation-pipeline:compiler}{compiler} checks every call against
them. To do so it must have seen the function's types before the call: a C file is read once, from top to bottom.

\begin{definition}[Function declaration, function definition]\label{def:b1:c-functions-and-stack:declaration}
A \emph{function declaration}\index{function declaration}, or \emph{prototype}\index{prototype}, gives a function's
name, the types of its parameters and the type of its result, followed by a semicolon:
\code{int digit_sum(unsigned n);}. A \emph{function definition}\index{function definition} gives the same heading
followed by the body between braces. A function has exactly one definition in a program, and may be declared many
times; every call must come after a declaration or the definition.
\end{definition}

A function whose result type is \code{void} returns nothing, and a parameter list \code{(void)} means ``no
parameters''. The header files of \cref{ch:b1:compilation-pipeline} are, for the most part, lists of
declarations: \texttt{<stdio.h>} declares \code{printf}, whose definition is in the C library. The order of
definitions in a file matters only through this rule; two functions that call each other need a declaration of
the second before the first (\cref{lst:b1:c-functions-and-stack:even-odd}).

\begin{definition}[Passing by value]\label{def:b1:c-functions-and-stack:by-value}
C passes every argument \emph{by value}\index{passing by value}: at a call, each argument expression is evaluated
and its value is copied into the corresponding parameter, which is a new variable of the called function. An
assignment to a parameter changes the copy, never the caller's variable.
\end{definition}

That is the whole story of \code{swap}: \code{a} and \code{b} were copies of \code{x} and \code{y}. A function hands a
value back to its caller through its result; to change a caller's variable it needs that variable's address, the
subject of \cref{ch:b1:pointers-and-arrays}. Arrays are the exception that is not one: \cref{ch:b1:pointers-and-arrays}
shows that the value copied for an array argument is its address.

\section{The call stack and its frames}\label{sec:b1:c-functions-and-stack:stack}

When \code{main} calls \code{swap}, the processor must jump to the code of \code{swap}, give it room for
\code{a}, \code{b} and \code{t}, and, when \code{swap} finishes, come back to the instruction after the call in
\code{main} with \code{main}'s variables intact. All of this is done with one region of memory, the \omterm{def:b1:c-functions-and-stack:call-stack}{call stack} met
in \cref{ch:b1:compilation-pipeline}.

\begin{definition}[Call stack, stack frame]\label{def:b1:c-functions-and-stack:call-stack}
The \emph{call stack}\index{call stack} is a region of memory used last in, first out: each call of a function
pushes a \emph{stack frame}\index{stack frame} (or activation record) holding what that call needs, and the return
pops it. A frame holds the call's parameters and local variables that do not fit in registers, the registers it
must save for its caller, and the \omterm{def:b1:c-functions-and-stack:return-address}{return address}. Only the frames of calls that have started and not yet returned
are on the stack, the most recent on top.
\end{definition}

\begin{definition}[Return address]\label{def:b1:c-functions-and-stack:return-address}
The \emph{return address}\index{return address} of a call is the address of the instruction that follows the call
instruction in the caller. The call instruction pushes it on the stack before jumping to the called function, and
the function's final \code{ret} instruction pops it and jumps there.
\end{definition}
% ledger: b1:psabi-call

On this machine the conventions are fixed by a document every \omterm{def:b1:compilation-pipeline:compiler}{compiler} for 64-bit x86 Linux follows, the System~V
application binary interface: the \code{call} instruction pushes the 8-byte \omterm{def:b1:c-functions-and-stack:return-address}{return address}, and at the entry of
every function the stack pointer, plus 8, is a multiple of 16. On x86-64 the stack grows towards lower addresses:
a new frame sits below its caller's. The program below prints, at each depth of a recursive call, how far below the
first call's local variable the current one lies.

\omcode{code/bachelor-1/03-c-functions-and-stack/listings/frames.c}{3}{23}{Measuring the distance between
frames}\label{lst:b1:c-functions-and-stack:frames}

\omoutput{code/bachelor-1/03-c-functions-and-stack/outputs/frames.txt}{Built at \code{-O0}: each call's frame is 64
bytes below the previous one.}\label{out:b1:c-functions-and-stack:frames}

Each frame of \code{descend} takes 64 bytes: the \omterm{def:b1:c-functions-and-stack:return-address}{return address} (8), the caller's saved frame pointer (8), and
48 bytes reserved below them, which hold the parameter and the two locals stored at \code{-O0} (16 bytes), an
8-byte guard value that this \omterm{def:b1:compilation-pipeline:compiler}{compiler} adds next to the locals of any function that takes a local's address, and
padding that keeps the 16-byte alignment. The printed differences
are the same at every run, though the addresses themselves are not; this is why the program prints differences.
At \code{-O2} the \omterm{def:b1:compilation-pipeline:compiler}{compiler} keeps most variables in registers, and the frames shrink or vanish: a measurement of
frames is a measurement of one build.

\begin{omfigure}
\begin{tikzpicture}[f/.style={draw, minimum width=46mm, minimum height=6mm, font=\small, anchor=north},
                    >={Stealth[length=2mm]}]
\node[f, fill=omCode!6] (m1) at (0,0) {\code{x = 1}, \code{y = 2}};
\node[f, fill=omCode!6] (m2) at (0,-0.6) {saved registers of \code{main}};
\node[f, fill=omThm!10] (r) at (0,-1.2) {return address (into \code{main})};
\node[f, fill=omDef!8] (s1) at (0,-1.8) {\code{a = 2}, \code{b = 1}};
\node[f, fill=omDef!8] (s2) at (0,-2.4) {\code{t = 1}};
\draw[decorate, decoration={brace, amplitude=4pt}] (2.45,0) -- (2.45,-1.2)
  node[midway, right=5pt, font=\small, align=left] {frame of \code{main}};
\draw[decorate, decoration={brace, amplitude=4pt}] (2.45,-1.2) -- (2.45,-3.0)
  node[midway, right=5pt, font=\small, align=left] {frame of \code{swap}};
\draw[->, thick] (-3.3,-3.0) node[left, font=\small] {stack pointer} -- (-2.35,-3.0);
\draw[->] (-3.6,0.2) -- node[left, font=\scriptsize, align=right] {lower\\addresses} (-3.6,-2.6);
\end{tikzpicture}
\par\smallskip{\small The stack while \code{swap} prints its line, the latest frame at the bottom. The parameters
\code{a} and \code{b} were filled with copies of \code{x} and \code{y}, then exchanged; \code{main}'s variables
are untouched. When \code{swap} returns, its frame is popped and the processor jumps to the \omterm{def:b1:c-functions-and-stack:return-address}{return address}.}
\end{omfigure}

The figure is a model. At \code{-O2}, \code{a}, \code{b} and \code{t} would live in registers and never reach the
stack; what is certain is the \omterm{def:b1:c-functions-and-stack:return-address}{return address}, and the rule that every call that has not returned keeps its frame.

\section{Recursion on the stack}\label{sec:b1:c-functions-and-stack:recursion}

A recursive call is a call like any other: it pushes a new frame, with new copies of the parameters and locals.
Here are three recursive functions on numbers; each stops at a base case and calls itself on a smaller number.

\omcode{code/bachelor-1/03-c-functions-and-stack/listings/digits.c}{2}{33}{Digit sum, binary writing and fast
power}\label{lst:b1:c-functions-and-stack:digits}

\omoutput{code/bachelor-1/03-c-functions-and-stack/outputs/digits.txt}{Three results.}

The call \code{digit_sum(2026)} pushes frames for 2026, 202, 20 and 2: four frames, one per digit, before the
first return. \code{print_binary} prints after its recursive call, so the leading bits come out first.
\code{power} halves its exponent at each call, so \code{power(3, 13)} makes five calls (exponents 13, 6, 3, 1, 0),
not thirteen. A function that calls itself twice builds a richer structure.

\begin{definition}[Tree of activations]\label{def:b1:c-functions-and-stack:call-tree}
The \emph{tree of activations}\index{tree of activations} (or \emph{call tree}\index{call tree}) of a run has one
node per call; the children of a node are the calls made directly by that call, in order. Its root is the first
call, its leaves the calls that made no call. The number of calls is the number of nodes; the greatest number of
frames on the stack at once is the height of the tree, counted in nodes.
\end{definition}

\omcode{code/bachelor-1/03-c-functions-and-stack/listings/calltree.c}{2}{20}{Printing the \omterm{def:b1:c-functions-and-stack:call-tree}{tree of activations},
one indented line per call}\label{lst:b1:c-functions-and-stack:calltree}

\omoutput{code/bachelor-1/03-c-functions-and-stack/outputs/calltree.txt}{Eleven calls, never more than four frames
at once.}\label{out:b1:c-functions-and-stack:calltree}

\begin{omfigure}
\begin{tikzpicture}[level distance=12mm, every node/.style={tnode, minimum size=9mm, font=\scriptsize\ttfamily},
  level 1/.style={sibling distance=52mm}, level 2/.style={sibling distance=26mm},
  level 3/.style={sibling distance=13mm}, leaf/.style={fill=omThm!12}, live/.style={very thick, draw=omDef}]
\node[live] {4,2}
  child {node {3,1}
    child {node[leaf] {2,0}}
    child {node {2,1} child {node[leaf] {1,0}} child {node[leaf] {1,1}}}}
  child {node[live] {3,2}
    child {node[live] {2,1} child {node[leaf] {1,0}} child {node[leaf, live] {1,1}}}
    child {node[leaf] {2,2}}};
\end{tikzpicture}
\par\smallskip{\small The \omterm{def:b1:c-functions-and-stack:call-tree}{tree of activations} of \code{binom(4, 2)}, nodes labelled $(n, k)$, leaves shaded. Thick
outlines: the four frames on the stack when the call \code{binom(1, 1)} of the right subtree runs, which is a path
from the root. The run makes 11 calls, the number of nodes, but needs at most 4 frames, the height.}
\end{omfigure}

Each leaf returns 1, and every result is a sum of leaf results, so the tree of \code{binom(n, k)} has
$\binom{n}{k}$ leaves; being a tree in which every inner node has two children, it has $2\binom{n}{k} - 1$ nodes in
all (\cref{exo:b1:c-functions-and-stack:8}). Calls grow like the result, which is huge; the depth is at most $n+1$.
That gap, between the work done and the memory held, is the theme of the weekend problem.

\begin{definition}[Mutual recursion]\label{def:b1:c-functions-and-stack:mutual}
Two or more functions are \emph{mutually recursive}\index{mutual recursion} when each calls another, so that a call
of one can lead, through the others, to a call of itself. In C the first one called before its definition must be
declared beforehand.
\end{definition}

\omcode{code/bachelor-1/03-c-functions-and-stack/listings/even_odd.c}{2}{22}{Even and odd, each defined through
the other}\label{lst:b1:c-functions-and-stack:even-odd}

\omoutput{code/bachelor-1/03-c-functions-and-stack/outputs/even_odd.txt}{Parity by \omterm{def:b1:c-functions-and-stack:mutual}{mutual
recursion}.}

Parity is a toy, but \omterm{def:b1:c-functions-and-stack:mutual}{mutual recursion} is the natural shape of anything defined by a grammar whose rules refer to
each other. An arithmetic expression is a sum of terms, a term is a product of factors, and a factor is a digit or
a whole expression in parentheses. One function per rule gives an evaluator in which \code{factor} calls
\code{expr}, which calls \code{term}, which calls \code{factor}.

\omcode{code/bachelor-1/03-c-functions-and-stack/listings/evaluator.c}{1}{36}{An evaluator in three \omterm{def:b1:c-functions-and-stack:mutual}{mutually
recursive} functions (the end of the file, \code{expr} and \code{evaluate}, follows the same
pattern)}\label{lst:b1:c-functions-and-stack:evaluator}

\omoutput{code/bachelor-1/03-c-functions-and-stack/outputs/calc.txt}{Products bind tighter than sums, as the
grammar says.}\label{out:b1:c-functions-and-stack:calc}

The depth of the recursion follows the nesting of the parentheses: \code{2*(3+(4*5))} reaches \code{factor} three
times on one path. The program that evaluates text from a user must therefore ask how deep a call chain can go.

\begin{definition}[Stack overflow]\label{def:b1:c-functions-and-stack:stack-overflow}
The \omterm{def:b1:c-functions-and-stack:call-stack}{call stack} has a maximal size, set when the program starts. A \emph{stack overflow}\index{stack overflow}
happens when a call needs a frame that does not fit in it any more. In C it is not reported by the language: the
program writes past the end of the stack region, and the system stops it.
\end{definition}

The program below makes each frame about 1 KiB, by giving the function a local array of 1000 bytes, and calls
itself until the stack runs out, printing each depth reached into a file. The shell command \code{ulimit -s},
covered with the shell in \cref{ch:b1:the-shell}, shows the stack limit of the programs it starts, in KiB, and
sets a smaller one; the parentheses keep the change to one command.

\omcode{code/bachelor-1/03-c-functions-and-stack/listings/deep.c}{2}{25}{Going down until the stack runs
out}\label{lst:b1:c-functions-and-stack:deep}

\omoutput{code/bachelor-1/03-c-functions-and-stack/outputs/depth.txt}{The program was stopped with exit status 139
at each limit; the last depth it reached grows with the limit.}\label{out:b1:c-functions-and-stack:depth}
% ledger: b1:linux-stack-default

The default limit on this system is 8192~KiB, that is 8~MiB. With 512~KiB and frames of 1040 bytes, the program
reached depth 500, and $512 \times 1024 / 1040 \approx 504.1$: the stack held the frames we predict, less about
two and a half kibibytes used before \code{dive} started (the program's arguments, its environment, the frames of
the start-up code and of \code{main}). The status 139 is how the system reports a program stopped for touching memory it may
not; \cref{ch:b1:the-shell} explains exit statuses, and the operating-system volume how the stop is delivered.
From one run to the next the last depth can differ by a few levels, because the system starts the stack at a
slightly random place, and from one machine to another because the environment differs; the run shown turned the
randomness off and started the program with an empty environment.

\begin{omfigure}
\begin{tikzpicture}
\begin{axis}[width=10cm, height=5.6cm, xlabel={stack limit $S$ (KiB)}, ylabel={depth reached},
  xmin=0, xmax=2200, ymin=0, ymax=2200, xtick={256,512,1024,2048}, ytick={0,500,1000,1500,2000},
  legend pos=north west, legend style={font=\small}, tick label style={font=\small}]
\addplot[domain=0:2150, thick, omDef, samples=2] {x*1024/1040};
\addlegendentry{$S / s$ with $s = 1040$ bytes}
\addplot[only marks, mark=*, omThm] coordinates {(256,248) (512,500) (1024,1004) (2048,2013)};
\addlegendentry{measured last depth}
\end{axis}
\end{tikzpicture}
\par\smallskip{\small The last depth reached by \cref{lst:b1:c-functions-and-stack:deep} for four stack limits.
The points lie 3 or 4 levels below the line $S/s$: the frames that fit in what is left once a
constant start-up cost, about 2.5~KiB, is paid. A loop doing the same work would use one frame at every limit.}
\end{omfigure}

When the last action of a function is a call whose result it returns unchanged, as \code{sum_to} does below,
that call is a tail call: nothing in the frame is needed after the inner call returns. An optimising \omterm{def:b1:compilation-pipeline:compiler}{compiler} may then
reuse the frame, turning the recursion into a jump. GCC does it at \code{-O2}, not at \code{-O0}; the language does
not require it, and the Year~2 volume treats the transformation. A program must not count on it.

\omcode{code/bachelor-1/03-c-functions-and-stack/listings/sum_to.c}{2}{15}{A tail-recursive
sum}\label{lst:b1:c-functions-and-stack:sum-to}

\omoutput{code/bachelor-1/03-c-functions-and-stack/outputs/tail.txt}{Ten million frames do not fit in 8~MiB; at
\code{-O2} there is only one.}\label{out:b1:c-functions-and-stack:tail}

\section{Scope and lifetime}\label{sec:b1:c-functions-and-stack:scope}

Two different questions are asked about a variable: where in the text its name can be used, and when, during the
run, the variable exists. They have different answers, and confusing them causes the bug at the end of this section.

\begin{definition}[Scope, shadowing]\label{def:b1:c-functions-and-stack:scope}
The \emph{scope}\index{scope} of a name is the part of the program text where it refers to a given declaration. A
name declared inside a block (between braces, or in the parentheses of a \code{for}) has \emph{block scope}: from
its declaration to the end of the block. A name declared outside every function has \emph{file scope}: from its
declaration to the end of the file. A declaration in an inner block \emph{shadows}\index{shadowing} a declaration of
the same name in an enclosing scope: inside the inner block, the name refers to the inner declaration.
\end{definition}

\omcode{code/bachelor-1/03-c-functions-and-stack/listings/shadow.c}{2}{21}{One name, three
variables}\label{lst:b1:c-functions-and-stack:shadow}

\omoutput{code/bachelor-1/03-c-functions-and-stack/outputs/shadow.txt}{Each use finds the nearest enclosing
declaration in the text.}

\code{show} sees the global \code{level}, not \code{main}'s: \omterm{def:b1:c-functions-and-stack:scope}{scope} is a property of the text, not of who called
whom. \omterm{def:b1:c-functions-and-stack:scope}{Shadowing} is legal and almost always a mistake waiting to happen; \code{-Wshadow} warns about it.

\begin{definition}[Lifetime; automatic, static and global variables]\label{def:b1:c-functions-and-stack:lifetime}
The \emph{lifetime}\index{lifetime (of a variable)} (or storage duration) of a variable is the part of the run
during which it exists and keeps its value. An \emph{automatic variable}\index{automatic variable}, a parameter or a
local declared without \code{static}, lives from the entry into its block to the exit from it, in the current
frame; each call has its own. A \emph{static variable}\index{static variable}, a local declared with
\code{static}, has block \omterm{def:b1:c-functions-and-stack:scope}{scope} but lives for the whole run, in the \omterm{def:b1:compilation-pipeline:layout}{data segment}, and is initialised once, before
\code{main} starts. A \emph{global variable}\index{global variable}, declared at file \omterm{def:b1:c-functions-and-stack:scope}{scope}, also lives for the
whole run, and its name is visible to the whole file.
\end{definition}

\omcode{code/bachelor-1/03-c-functions-and-stack/listings/counter_static.c}{2}{19}{A static local counts the
calls}\label{lst:b1:c-functions-and-stack:counter}

\omoutput{code/bachelor-1/03-c-functions-and-stack/outputs/counter_static.txt}{\code{visits} survives the calls,
\code{fresh} does not.}

\begin{definition}[Internal linkage]\label{def:b1:c-functions-and-stack:linkage}
A function or a variable declared at file \omterm{def:b1:c-functions-and-stack:scope}{scope} with \code{static} has \emph{internal
linkage}\index{internal linkage}: its name is local to its file, and the \omterm{def:b1:compilation-pipeline:compiler}{linker} never matches it with a name of
another file. Without \code{static}, a \omterm{def:b1:c-functions-and-stack:scope}{file-scope} name has \emph{external linkage}: it is visible to the \omterm{def:b1:compilation-pipeline:compiler}{linker}, as
the global symbols of \cref{ch:b1:compilation-pipeline}.
\end{definition}

The word \code{static} thus means two things, depending on where it is written: inside a function it changes the
\omterm{def:b1:c-functions-and-stack:lifetime}{lifetime}; at file \omterm{def:b1:c-functions-and-stack:scope}{scope} it changes the linkage. In the evaluator, \code{p}, \code{good}, \code{factor},
\code{term} and \code{expr} are \code{static}: helpers that no other file needs, which cannot then collide with
another file's \code{expr}. A module shows the rest of the program only what it declares in its header, a practice
\cref{ch:b1:c-modules-make-tests} makes systematic.

\begin{omfigure}
\begin{tikzpicture}[>={Stealth[length=2mm]}, bar/.style={line width=5pt}]
\draw[->] (0,0) -- (11,0) node[right, font=\small] {time};
\foreach \x/\t in {0.4/start, 2/call 1, 4.2/return, 6/call 2, 8.2/return, 10.4/end}
  {\draw (\x,0.1) -- (\x,-0.1) node[below, font=\scriptsize] {\t};}
\draw[bar, omCode] (0.4,4) -- (10.4,4); \node[left, font=\small] at (0.3,4) {global \code{level}};
\draw[bar, omCode] (0.4,3.2) -- (10.4,3.2); \node[left, font=\small] at (0.3,3.2) {static \code{visits}};
\draw[bar, omDef] (2,2.4) -- (4.2,2.4); \draw[bar, omDef] (6,2.4) -- (8.2,2.4);
\node[left, font=\small] at (0.3,2.4) {automatic \code{fresh}};
\draw[bar, omDef] (2.4,1.6) -- (3.8,1.6); \node[left, font=\small] at (0.3,1.6) {inner block's local};
\draw[bar, omThm] (4.2,0.8) -- (10.4,0.8); \node[left, font=\small] at (0.3,0.8) {returned address};
\node[font=\scriptsize, omThm, above] at (7.3,0.9) {points to a dead variable};
\end{tikzpicture}
\par\smallskip{\small \omterm{def:b1:c-functions-and-stack:lifetime}{Lifetimes} along a run. Global and \omterm{def:b1:c-functions-and-stack:lifetime}{static variables} exist from start to end. An \omterm{def:b1:c-functions-and-stack:lifetime}{automatic
variable} exists once per call (two calls, two separate lives), a block's local only inside its block. An address
returned from a call outlives the variable it names.}
\end{omfigure}

The last line of the figure is a classic bug. A function that returns the address of one of its \omterm{def:b1:c-functions-and-stack:lifetime}{automatic
variables} returns the address of something that ceases to exist at the return; its frame will be reused by the
next call.

\omcode{code/bachelor-1/03-c-functions-and-stack/listings/dangling_local.c}{2}{15}{Returning the address of a
local}\label{lst:b1:c-functions-and-stack:dangling}

\omoutput{code/bachelor-1/03-c-functions-and-stack/outputs/dangling_local.txt}{GCC warns, then returns a null
address.}\label{out:b1:c-functions-and-stack:dangling}
% ledger: b1:gcc-return-local-addr

Using the dead variable is \omterm{def:b1:c-types-and-control:ub}{undefined behaviour}, and GCC, knowing this, chose to return the null address, which no
object has: a later use crashes at once instead of reading garbage. Another \omterm{def:b1:compilation-pipeline:compiler}{compiler} may return the stale address.
The fix is to return a value, or an object whose \omterm{def:b1:c-functions-and-stack:lifetime}{lifetime} the caller controls: a \omterm{def:b1:c-functions-and-stack:lifetime}{static variable}, a variable of the
caller passed by address (\cref{ch:b1:pointers-and-arrays}), or a block of the heap
(\cref{ch:b1:dynamic-memory}).

\begin{history}[The stack before C]
Early languages such as Fortran gave each function fixed places for its variables, which ruled out recursion: a
second call would overwrite the first one's values. Algol~60 required recursion, and with it a stack of activation
records; its implementers made the stack the standard way to compile procedure calls, and every language since,
C included, has inherited it.
\end{history}

\section{Exercises}\label{sec:b1:c-functions-and-stack:exercises}

\begin{exercise}[$\star$]\label{exo:b1:c-functions-and-stack:1}
Predict the line printed by the program below, and explain the \omterm{def:b1:compilation-pipeline:compiler}{compiler}'s warning.
\omcode{code/bachelor-1/03-c-functions-and-stack/listings/exo_reset.c}{2}{22}{Two resets}
\end{exercise}

\begin{exercise}[$\star$]\label{exo:b1:c-functions-and-stack:2}
Predict the two lines printed, naming for each use of \code{x} the declaration it refers to.
\omcode{code/bachelor-1/03-c-functions-and-stack/listings/exo_shadow.c}{2}{24}{Which \code{x}?}
\end{exercise}

\begin{exercise}[$\star$]\label{exo:b1:c-functions-and-stack:3}
Predict the three numbers printed.
\omcode{code/bachelor-1/03-c-functions-and-stack/listings/exo_static.c}{2}{19}{A total kept between calls}
\end{exercise}

\begin{exercise}[$\star$]\label{exo:b1:c-functions-and-stack:4}
Draw the \omterm{def:b1:c-functions-and-stack:call-stack}{call stack} at the moment the marked line runs for the second time: which frames are present, in which
order, and what values do their parameters and locals hold? Which \omterm{def:b1:c-functions-and-stack:return-address}{return address} does the top frame hold?
\omcode{code/bachelor-1/03-c-functions-and-stack/listings/exo_frames.c}{2}{18}{Frames at a marked line}
\end{exercise}

\begin{exercise}[$\star\star$]\label{exo:b1:c-functions-and-stack:5}
The program below prints \code{greeting: (null)}. Explain why, using the warning, and give two correct versions of
\code{greeting}.
\omcode{code/bachelor-1/03-c-functions-and-stack/listings/exo_name.c}{2}{16}{A greeting built in a local array}
\omoutput{code/bachelor-1/03-c-functions-and-stack/outputs/exo_name.txt}{The warning and the result.}
\end{exercise}

\begin{exercise}[$\star\star$]\label{exo:b1:c-functions-and-stack:6}
The function \code{steps} counts the steps from \code{n} down to zero, by twos. What happens for \code{steps(7)}?
Describe the chain of calls, say what the program does at \code{-O0}, and correct the base case so that
\code{steps} is defined for every \code{n} $\geq 0$.
\omcode{code/bachelor-1/03-c-functions-and-stack/listings/exo_steps.c}{2}{15}{Down by twos}
\end{exercise}

\begin{exercise}[$\star\star$]\label{exo:b1:c-functions-and-stack:7}
Write \code{int digit_sum_iter(unsigned long long n)}, a version of \code{digit_sum} with a loop instead of
recursion, and state a loop invariant that shows it computes the same value. How many frames does it use for a
20-digit number?
\end{exercise}

\begin{exercise}[$\star\star$]\label{exo:b1:c-functions-and-stack:8}
Prove that the naive \code{binom(n, k)} of \cref{lst:b1:c-functions-and-stack:calltree} makes exactly
$2\binom{n}{k} - 1$ calls, for $0 \leq k \leq n$. Deduce the number of calls for \code{binom(30, 15)}, given that
$\binom{30}{15} = 155\,117\,520$.
\end{exercise}

\begin{exercise}[$\star\star$]\label{exo:b1:c-functions-and-stack:9}
Trace the calls made by \code{is_odd(7)} in the program below and give its output. What would \code{is_odd(-1)} do?
\omcode{code/bachelor-1/03-c-functions-and-stack/listings/exo_parity.c}{2}{24}{Parity with a trace}
\end{exercise}

\begin{exercise}[$\star\star\star$]\label{exo:b1:c-functions-and-stack:10}
The naive Fibonacci function returns $n$ for $n < 2$ and \code{fib(n-1) + fib(n-2)} otherwise, with $F(0) = 0$,
$F(1) = 1$. Prove that \code{fib(n)}, for $n \geq 1$, makes $C(n) = 2F(n+1) - 1$ calls, and that at most $n$ of
them are on the stack at once. Compare the two for $n = 30$ ($F(31) = 1\,346\,269$).
\end{exercise}

\begin{exercise}[$\star\star\star$]\label{exo:b1:c-functions-and-stack:11}
Write a recursive function that solves the towers of Hanoi for $n$ discs, records each move (\code{A->C}), returns
the number of moves and the greatest depth of recursion reached. Prove that it makes $2^n - 1$ moves and reaches
depth $n + 1$ when the first call has depth 1 and the calls with no disc count.
\end{exercise}

\begin{exercise}[$\star\star\star$]\label{exo:b1:c-functions-and-stack:12}
A recursive function has frames of $s$ bytes, and the stack holds $S$ bytes, of which $c$ are used before the first
call. Give the greatest depth $d$ that fits, and prove it. Using the measurements of
\cref{out:b1:c-functions-and-stack:depth}, find $c$, and predict the last depth reached under a limit of 4096~KiB.
\end{exercise}

\section{Problem: How deep can a program go?}\label{sec:b1:c-functions-and-stack:problem}

\begin{problem}[{Weekend problem --- How deep can a program go?}]\label{pb:b1:c-functions-and-stack:1}
A recursive program uses two resources: time, counted in calls, and memory, counted in frames on the stack at once.
The problem measures both for several shapes of recursion and turns the stack limit into a predicted depth. It uses
the outputs of \cref{lst:b1:c-functions-and-stack:frames,lst:b1:c-functions-and-stack:deep}, and the two programs
below.
\omcode{code/bachelor-1/03-c-functions-and-stack/listings/shapes.c}{15}{31}{Three shapes, instrumented by a call to
\code{enter(depth)} (the rest of the file prints the table)}\label{lst:b1:c-functions-and-stack:shapes}
\omoutput{code/bachelor-1/03-c-functions-and-stack/outputs/shapes.txt}{Calls and greatest depth for $n = 5, 10, 15,
20$.}\label{out:b1:c-functions-and-stack:shapes}
\omcode{code/bachelor-1/03-c-functions-and-stack/listings/parity_deep.c}{7}{20}{\omterm{def:b1:c-functions-and-stack:mutual}{Mutual recursion} ten million
deep}
\omoutput{code/bachelor-1/03-c-functions-and-stack/outputs/problem.txt}{The same program at \code{-O0} and
\code{-O2}, and the calls left in the \code{-O2} assembly.}\label{out:b1:c-functions-and-stack:problem}

\medskip\noindent\textbf{Part I --- Frames measured.}
\begin{enumerate}
\item In \cref{out:b1:c-functions-and-stack:frames}, what is the size of one frame of \code{descend}? Why do we
print differences of addresses rather than addresses?
\item List what a frame of \code{descend} must contain at \code{-O0}, with sizes, and check that it can fit in that
size with the alignment rule (the stack pointer plus 8 is a multiple of 16 at each function's entry).
\item In \cref{out:b1:c-functions-and-stack:depth}, the frame of \code{dive} measures 1040 bytes. Why more than
1000? Why is 1040 a multiple of 16?
\item The program measures the frame size at depth 1, as the distance from depth 0. Why would it be wrong to
measure it as the distance from a local of \code{main}?
\item Why may the same programs give smaller frames at \code{-O2}? Which part of a frame can never disappear for a
call that is not a tail call?
\item Why must \code{deep.c} call \code{fflush} at each level, and why is its output sent to a file?
\end{enumerate}

\medskip\noindent\textbf{Part II --- Calls against depth.}
\begin{enumerate}[resume]
\item Prove that \code{count_down(n, 1)} makes $n + 1$ calls and reaches depth $n + 1$; check the table.
\item Prove by induction that \code{hanoi(n, 1)} makes $2^{n+1} - 1$ calls and reaches depth $n + 1$; check the
table.
\item For \code{fib}, the table gives 15 calls for $n = 5$. Check it against $2F(n+1) - 1$, with $F(6) = 8$.
\item Why does \code{fib(n, 1)} reach depth $n$, not $n + 1$?
\item Between $n = 15$ and $n = 20$, by what factor did the calls of \code{fib} grow? Compare with
$\varphi^5$, where $\varphi = (1 + \sqrt 5)/2$.
\item For each shape, give the ratio of calls to depth at $n = 20$.
\item The \omterm{def:b1:c-functions-and-stack:call-tree}{tree of activations} of \code{hanoi(n, 1)} is a complete binary tree. Which nodes are its leaves, and how
many are there?
\end{enumerate}

\medskip\noindent\textbf{Part III --- Mutual recursion and tail calls.}
\begin{enumerate}[resume]
\item How many calls does \code{is_even(10000000)} make, and what is the greatest number of frames on the stack at
\code{-O0}? Why did the \code{-O0} program stop with status 139?
\item In the \code{-O2} assembly, \code{is_odd} ends with \code{jmp is_even} instead of \code{call is_even}. What
does that change on the stack?
\item The only \code{call} left is \code{main}'s. Why can the \omterm{def:b1:compilation-pipeline:compiler}{compiler} not replace it by a jump?
\item Is the call \code{fib(n - 1, depth + 1)} in \code{fib} a tail call? And the call
\code{count_down(n - 1, depth + 1)} in \code{1 + count_down(...)}?
\item Rewrite \code{count_down} with an accumulator so that its recursive call is a tail call.
\end{enumerate}

\medskip\noindent\textbf{Part IV --- The stack budget.}
\begin{enumerate}[resume]
\item With $S = 512$~KiB and $s = 1040$ bytes, compute $\lfloor S / s \rfloor$ and compare with the last depth
reached, 500.
\item Do the same for the three other limits of \cref{out:b1:c-functions-and-stack:depth}. Is the difference
always the same? Explain why a constant cost counted in bytes, not in levels, accounts for all four.
\item What does that cost consist of? Estimate in bytes the stack used before \code{dive(0)} starts.
\item Under the default limit, 8192~KiB, predict the last depth reached by \code{deep.c}.
\item Under the default limit and at \code{-O0}, how deep can \code{descend} of \cref{lst:b1:c-functions-and-stack:frames}
go, roughly?
\item The evaluator nests one call of \code{expr}, \code{term} and \code{factor} per pair of parentheses. If its
three frames take 48 bytes each, how many nested parentheses can a text have under 8192~KiB? What should a program
that evaluates untrusted text do about it?
\item State the named result of the problem, \emph{the depth bound}, with the values measured here.
\end{enumerate}
\end{problem}
